\begin{proposition}\label{prop:decic-no56}
An elliptic curve over a number field of degree $10$ has no point
of order $56$.
\end{proposition}
\begin{proof}
Put $X=X_1(56)$ and $Y=X_1(2,28)$. Both Jacobians have rational
rank zero by \cite[Theorem~3.1]{DEHMZ}. We have
\[
 g(X)=61,\qquad g(Y)=25,\qquad
 \operatorname{gon}_{\mathbb Q}X\geq\frac{7\cdot1152}{800}>10,
 \qquad \operatorname{gon}_{\mathbb Q}Y\geq12,
\]
where the gonality bounds follow from \cite{Abramovich} and
\cite[Table~1]{NV}, respectively.
Suppose that $X$ has a noncuspidal effective rational divisor $D$
of degree $10$. We use reduction at $3$ throughout the proof.

First, $\overline D$ cannot be entirely cuspidal. Indeed, the
cuspidal boundary is finite \'etale at $3$, so it has a cuspidal
lift $C$ of degree $10$ over $\mathbb Q_3$. The class $[D-C]$ is
torsion by rank zero and Manin--Drinfeld, and reduces to zero.
The formal kernel over $\mathbb Q_3$ is torsion-free. Thus $D\sim C$,
so the global degree-$10$ line bundle has at least two sections
after base change to $\mathbb Q_3$, and hence over $\mathbb Q$.
This contradicts the gonality bound.

We first apply necessary class tests on $Z=X_1(28)$.
The map $X\to Y$ sends $(E,R)$ to
\[
 (E/\langle28R\rangle,P_2,\phi(R)),
\]
where $P_2$ is the nonzero point in the kernel of the dual isogeny.
The point $P_2$ is independent of $14\phi(R)$, since the dual sends
the latter to $28R\ne0$.
We use the following four maps from $Y$ to $Z$:
\[
 (E,P_2,Q)\longmapsto(E,Q),\quad(E,P_2+Q),\quad
 (E/\langle P_2\rangle,\phi Q),\quad
 (E/\langle P_2+14Q\rangle,\psi Q).
\]
The complete rational torsion subgroup of $J_1(28)$ is cuspidal
and isomorphic to $[2,4,12,936]$ by
\cite[Theorem~4.13]{DEHMZ}. Its reduction at $3$ has index $4$
in $J_1(28)(\mathbb F_3)\simeq[4,4,24,936]$. Each of the four
image classes must belong to this rational image, with its
cuspidal complement and all residue-field norms retained.

The numbers of noncuspidal closed points on $X_{\mathbb F_3}$
in degrees $1,\ldots,10$ are
\[
             0,0,0,0,72,48,324,588,2112,6300.
\]
The cuspidal degrees have signature $1^{18}2^9 3^2 6^5$.
The complete source census has $797$ diamond orbits. There are
$1229$ support patterns to test: one noncuspidal orbit with a
cuspidal complement, or a sum of two degree-$5$ places, allowing
repetition. Exactly $135$ pass the four necessary class tests.
Their noncuspidal degrees and numbers are
\[
\begin{array}{c|rrrrrr}
 \text{degrees}&6&7&8&9&10&5+5\\\hline
 \text{number}&2&6&11&8&48&60.
\end{array}
\]
The finite Picard presentation is proved. An independent native
replay verifies all $2560$ image classes, including the $2490$
nontrivial principal-divisor witnesses in both finite and infinite
orders. An independent marked V\'elu calculation verifies every
image place and norm, and an integer replay verifies all $1229$ tests.

We next apply the operator
$A_5=T_5-1-5\langle5\rangle$ on $Y$. It annihilates
$\operatorname{Jac}(Y)(\mathbb Q)$: this follows from
Eichler--Shimura, rank zero and injectivity of rational torsion
at the good odd prime $5$. It also annihilates a rational base
cusp as a divisor. Hence $A_5\overline E$ must be principal for
the projected degree-$10$ divisor $E$.
The $135$ source patterns give $41$ distinct marked patterns on
$Y$, up to the rational automorphisms
\[
 (P_2,Q)\longmapsto(P_2+14bQ,cP_2+aQ),\qquad
 b,c\in\{0,1\},\quad a\in\{1,3,5,9,11,13\}.
\]
For completeness, the cusp signature on $Y_{\mathbb F_3}$ is
$1^{12}2^6 3^4 6^2$. The four cubic cusps form two pairs of
global degree-$6$ cusps. All cuspidal complements have degree
at most $4$. The closed cusps of degrees $1$ and $2$ are killed
literally by $A_5$; allowing a cubic cusp gives the seventeen
possible corrections recorded in the shared certificate.
Testing all these corrections leaves the six marked patterns
\[
                       8,14,26,27,36,41.
\]
These labels refer to \texttt{Y228\_shared/union\_cases.m}.
They come from $15$ original source patterns, or $11$ source
diamond orbits. The complete remaining source list is
\[
\begin{array}{c|c|c}
 \text{source labels}&\text{noncuspidal degrees}&\text{image label on }Y\\\hline
 65&8&8\\
 271,327&10&14\\
 491,578&10&26\\
 492,737&10&27\\
 818,837&5+5&36\\
 978,1009&5+5&41.
\end{array}
\]
All normalizations act on the whole marked divisor. In particular,
the relative markings in a sum of two degree-$5$ places are retained.

We first exclude image label $27$ by a global calculation.
Let $K_0=\mathbb Q(a)$, where
\[
 a^5-a^4-2a^3-a^2+2a+2=0,
\]
and put
\[
 \begin{split}
 b&=1032-239a-222a^2-759a^3+419a^4,\\
 c&=42-11a-8a^2-31a^3+17a^4.
 \end{split}
\]
The curve $E:y^2+(1-c)xy-by=x^3-bx^2$ has a point
$P=(0,0)$ of order $28$. Adjoining an independent point $S$ of
order $2$ gives a quadratic extension $K/K_0$ and a marked point
of $Y(K)$. Its trace reduces to one degree-$10$ place, and exact
finite-field arithmetic identifies its automorphism orbit with
label $27$. An independent Sage calculation verifies this binding.

Any proposed global degree-$10$ divisor on $Y$ with this reduction
is linearly equivalent to the known trace: their difference is
rational torsion and reduces to zero at $3$. Since
$\operatorname{gon}_{\mathbb Q}Y\geq12$, the divisors are equal.
Thus the curve, field and marking are the displayed ones, up to
a rational automorphism of $Y$.
For this point to lift to $X$, its order-$28$ point $Q$ must lie
in the image of the dual $2$-isogeny. After translating its kernel
point $P_2$ to $(0,0)$ and completing the square, the dual has
first coordinate $(y/(2x))^2$. Consequently
\[
                         x(Q)-x(P_2)\in K^{\times2}
\]
is necessary. The standard $2$-descent map is a homomorphism, so
all $24$ possible markings reduce to the four choices
\[
 P_2=S+14\epsilon P,\qquad Q=P+\eta P_2,
 \qquad\epsilon,\eta\in\{0,1\}.
\]
All four squareclasses are nontrivial. Indeed, at good unramified
primes above $43$, with $a\mapsto37$ and $x(S)\mapsto13$ or $22$,
the four nonzero $x$-differences reduce respectively to
\[
                         30,39,30,27\in\mathbb F_{43}.
\]
Each has Legendre symbol $-1$. These residue certificates agree
with separate exact number-field square tests. This excludes
source labels $492$ and $737$.

It remains to test nine source patterns. We give an explicit
quadratic model which makes these calculations small.
Write the universal curve on $Y$ as
\[
 E':y^2=x\bigl(x^2+Cx+B(1-B-C)\bigr),\qquad
 P_2=(0,0),\quad Q=(B,B).
\]
Put $\Delta=C^2-4B(1-B-C)$. The function field of $X$ over that
of $Y$ is
\[
                           z^2=B.
\]
Indeed, on $E'':y^2=x(x^2-2Cx+\Delta)$ the point
\[
 R=(2B+C+2z,\ 2z(2B+C+2z))
\]
has order $56$. The dual isogeny
\[
 (x,y)\longmapsto
 \left(\frac{x+\Delta/x-2C}{4},
       \frac{y(1-\Delta/x^2)}8\right)
\]
sends $R$ to $Q$. Conversely, a lift of $Q$ determines the
square root $z$ and hence this marked curve. These identities
identify the quadratic cover with $X_1(56)$. They are verified
in the function field; its genus is $61$ and the odd part of
$\operatorname{div}(B)$ has degree $24$.

We use the following character decomposition in characteristic
$3$. Let $F=\mathbb F_3(Y)$, $L=F(z)$ and $z^2=B$. For a
source divisor $G$ of degree zero supported where the cover is
unramified, write $n_q$ for its coefficient at $q$. For each
base place $p$, set
\[
 a_p=\max\bigl(\{0\}\cup\{n_q:q\mid p\}\bigr),\qquad
 H=\sum_p a_pp,\qquad
 R_B=\left\lfloor\tfrac12\operatorname{div}(B)\right\rfloor.
\]
Then $G\leq\pi^*H$ and
\[
 \mathcal L_L(\pi^*H)
 =\mathcal L_F(H)\oplus z\mathcal L_F(H+R_B).
\]
To obtain $\mathcal L_L(G)$, it is enough to impose at every
$q\mid p$ the first $a_p-n_q$ vanishing conditions after
multiplication by a local parameter to the power $a_p$.
We impose these conditions on both sheets, including places
with $n_q=0$. If a place is inert, its full residue field is
used instead. Thus no conditions are lost by passing to the
quadratic model.

For the eight purely noncuspidal source patterns, take
$G=A_5\overline D$ on $X$. For source label $65$, write
$\overline D=P_8+C_2$ and use
$A_{11}=T_{11}-1-11\langle11\rangle$.
On Tate cusp labels, let $\sigma_q$ denote the cyclotomic Galois
action. The point-level formula
\[
       T_q=\sigma_q+q\langle q\rangle\sigma_q^{-1}
\]
shows that $A_{11}$ kills every closed cusp of degree at most
$2$ over $\mathbb F_3$, as also checked on all $27$ such cusps.
Thus in this case $G=A_{11}P_8$. Both operators annihilate the
rational source Jacobian and a rational base cusp, so the
displayed degree-zero divisor must be principal.

The exact computations give the following table. The last
column is the rank of the vanishing-condition matrix over
$\mathbb F_3$.
\[
\begin{array}{c|r|r|r|r|r}
 \text{source label}&q&\deg H&\ell(H)&\ell(H+R_B)&\text{rank}\\\hline
 271,327,491,578,818,837,978,1009&5&60&36&24&60\\
 65&11&96&72&60&132.
\end{array}
\]
There are respectively $120$ and $192$ scalar vanishing
conditions. In every case the matrix has full row rank, so
$\ell(G)=0$, a contradiction.

The certificate records the exact function-field model, the
base prime generators, every marked source and Hecke image,
and all residue norms. The Taylor expansions are obtained
from the plane equation by Hensel lifting, and the square root
is selected by the actual marked point of order $56$.
Independent native base-space and Sage series calculations
verify the bases, local equations and restriction matrices;
independent marked isogenies verify all six $5$-kernels and all
twelve $11$-kernels used in the corresponding tests.
This completes the proof.
\end{proof}
